% sec_circle_dimension_algebra.tex -- circle dimension and phase sampling
\section{Classical Phase Sampling and Circle Dimension}%
\label{sec:circle-dimension}

\subsection{Finite phase sampling}

The Cayley chart isolates the circle as the basic visible phase factor.
The algebraic material in this section is a finite sampling interface for
that phase coordinate, not a separate source of analytic novelty.  The
statements below are classical consequences of the structure theorem for
finitely generated abelian groups and Pontryagin duality
\cite{Hungerford1974,Folland1995}.  We include the proofs only to fix the
normalization of the phase-sampling counts used in
Proposition~\ref{prop:phase-coordinate-interface},
Proposition~\ref{prop:finite-grid-coefficient-recovery}, and the
finite phase-grid verification
Proposition~\ref{prop:eighth-order-finite-phase-verification}.

\begin{definition}[Circle dimension]\label{def:circle-dimension}
Let $G$ be a finitely generated abelian group and write
\[
\widehat G:=\operatorname{Hom}(G,\TT)
\]
for its Pontryagin dual, endowed with the compact-open topology.  We
define the \emph{circle dimension} of~$G$ by
\[
\cdim(G):=\dim\bigl((\widehat G)^0\bigr),
\]
where $(\widehat G)^0$ denotes the identity component of~$\widehat G$.
\end{definition}

\begin{definition}[Half-circle dimension]\label{def:half-circle-dimension}
Let $M$ be a finitely generated cancellative commutative monoid, and let
$G(M)$ be its Grothendieck group completion.  The \emph{half-circle
dimension} of~$M$ is
\[
\cdim_+(M):=\frac{1}{2}\cdim\bigl(G(M)\bigr).
\]
\end{definition}

\begin{proposition}[Basic computation]\label{prop:circle-dimension-basic}
If $G\simeq \ZZ^r\oplus T$ with $T$ finite, then
\[
\widehat G\simeq \TT^r\times \widehat T
\qquad\text{and}\qquad
\cdim(G)=r=\dim_{\QQ}(G\otimes_{\ZZ}\QQ).
\]
In particular,
\[
\cdim(\ZZ^r)=r,
\qquad
\cdim(T)=0
\qquad\text{for every finite abelian group }T,
\]
and
\[
\cdim_+(\NN^d)=\frac{d}{2}
\qquad(d\ge 1).
\]
\end{proposition}

\begin{proof}
Pontryagin duality sends direct sums to direct products, so
\[
\widehat G\simeq \widehat{\ZZ^r}\times \widehat T \simeq \TT^r\times \widehat T.
\]
Since $\widehat T$ is finite, its identity component is trivial, hence
$(\widehat G)^0\simeq \TT^r$ and $\cdim(G)=r$.

The tensor product $G\otimes_{\ZZ}\QQ$ kills the torsion subgroup~$T$,
so
\[
G\otimes_{\ZZ}\QQ \simeq \QQ^r,
\]
and therefore $\dim_{\QQ}(G\otimes_{\ZZ}\QQ)=r$.  The special cases for
$\ZZ^r$ and finite groups are immediate.  Finally,
$G(\NN^d)\simeq \ZZ^d$, so
\[
\cdim_+(\NN^d)=\frac{1}{2}\cdim(\ZZ^d)=\frac{d}{2}.
\]
\end{proof}

\begin{theorem}[Classical uniqueness of the rank recovery]%
\label{thm:cdim-axiomatic-rigidity}
Let $f$ be a function from the class of finitely generated abelian groups
to $\NN\cup\{0\}$ satisfying:
\begin{enumerate}[label=(\roman*)]
\item $f(G)=f(H)$ whenever $G\simeq H$;
\item $f(G\oplus H)=f(G)+f(H)$ for all $G,H$;
\item if
      $0\to F\to G\to H\to 0$
      is exact with $F$ finite, then $f(G)=f(H)$;
\item $f(\ZZ)=1$;
\item $f(F)=0$ for every finite abelian group~$F$.
\end{enumerate}
Then $f(G)=\cdim(G)$ for every finitely generated abelian group~$G$.
\end{theorem}

\begin{proof}
By the structure theorem for finitely generated abelian groups
\cite{Hungerford1974}, every finitely generated abelian group admits a
decomposition
\[
G\simeq \ZZ^r\oplus T
\]
with $T$ finite.  By (iii) and (v),
\[
f(G)=f(\ZZ^r).
\]
Using (ii) and (iv) gives $f(\ZZ^r)=r$.  Proposition~\ref{prop:circle-dimension-basic}
shows that $\cdim(G)=r$, hence $f(G)=\cdim(G)$.
\end{proof}

\begin{theorem}[Short exact additivity and rank-nullity]%
\label{thm:cdim-short-exact-additivity}
Let
\[
0\longrightarrow A\longrightarrow B\longrightarrow C\longrightarrow 0
\]
be a short exact sequence of finitely generated abelian groups.  Then
\[
\cdim(B)=\cdim(A)+\cdim(C).
\]
Consequently, every homomorphism $\varphi:B\to D$ satisfies
\[
\cdim(B)=\cdim(\ker\varphi)+\cdim(\operatorname{im}\varphi).
\]
\end{theorem}

\begin{proof}
By Proposition~\ref{prop:circle-dimension-basic},
\[
\cdim(G)=\dim_{\QQ}(G\otimes_{\ZZ}\QQ)
\]
for every finitely generated abelian group~$G$.  Since $\QQ$ is a flat
$\ZZ$-module, tensoring the given short exact sequence by~$\QQ$ yields
another short exact sequence
\[
0\longrightarrow A\otimes_{\ZZ}\QQ
\longrightarrow B\otimes_{\ZZ}\QQ
\longrightarrow C\otimes_{\ZZ}\QQ
\longrightarrow 0.
\]
Taking dimensions gives the first identity.  The second follows by
applying the first to
\[
0\longrightarrow \ker\varphi
\longrightarrow B
\longrightarrow \operatorname{im}\varphi
\longrightarrow 0.
\]
\end{proof}

\begin{proposition}[Tensor, Hom, and Ext laws]%
\label{prop:cdim-tensor-hom-ext-laws}
Let $G$ and $H$ be finitely generated abelian groups.  Then:
\begin{enumerate}[label=(\roman*)]
\item
      $\cdim(G\otimes_{\ZZ}H)=\cdim(G)\cdot \cdim(H)$;
\item
      $\cdim(\operatorname{Hom}(G,H))=\cdim(G)\cdot \cdim(H)$;
\item
      $\cdim(\operatorname{Ext}^{1}(G,H))=0$.
\end{enumerate}
\end{proposition}

\begin{proof}
Write
\[
G\simeq \ZZ^r\oplus T,
\qquad
H\simeq \ZZ^s\oplus E,
\]
with $T$ and $E$ finite.  By
Proposition~\ref{prop:circle-dimension-basic},
$r=\cdim(G)$ and $s=\cdim(H)$.

For (i), the free summand of $G\otimes_{\ZZ}H$ is
\[
\ZZ^r\otimes_{\ZZ}\ZZ^s\simeq \ZZ^{rs},
\]
while all other tensor products are torsion.  Hence
$\cdim(G\otimes_{\ZZ}H)=rs$.

For (ii), one has
\[
\operatorname{Hom}(G,H)\simeq
\operatorname{Hom}(\ZZ^r,H)\oplus \operatorname{Hom}(T,H).
\]
The first summand is isomorphic to $H^r$, whose free rank is $rs$; the
second is finite because $T$ is finite and $H$ is finitely generated.
Thus $\cdim(\operatorname{Hom}(G,H))=rs$.

For (iii), additivity in the first variable gives
\[
\operatorname{Ext}^{1}(G,H)\simeq
\operatorname{Ext}^{1}(\ZZ^r,H)\oplus \operatorname{Ext}^{1}(T,H).
\]
The first term vanishes, while $T$ is a finite direct sum of cyclic
groups $C_n=\ZZ/n\ZZ$, and
\[
\operatorname{Ext}^{1}(C_n,H)\simeq H/nH.
\]
Since $H$ is finitely generated, each quotient $H/nH$ is finite, so
$\operatorname{Ext}^{1}(T,H)$ is finite.  Therefore
$\cdim(\operatorname{Ext}^{1}(G,H))=0$.
\end{proof}

\subsection{Phase sampling and reconstruction}

\begin{definition}[Phase-sampling counts]\label{def:cdim-phase-spectrum}
For a finitely generated abelian group $G$ and an integer $N\ge 1$, let
\[
\mu_N:=\{z\in \TT:\ z^N=1\}.
\]
We define the \emph{phase-sampling count} of~$G$ at level~$N$ by
\[
\Sigma_G(N):=\bigl|\operatorname{Hom}(G,\mu_N)\bigr|.
\]
For $N\ge 2$ we also set
\[
\delta_G(N):=\frac{\log \Sigma_G(N)}{\log N}.
\]
\end{definition}

\begin{proposition}[Sampling counts and finite quotients]%
\label{prop:cdim-phase-spectrum-quotient}
For every finitely generated abelian group~$G$ and every integer
$N\ge 1$,
\[
\Sigma_G(N)=|G/NG|.
\]
\end{proposition}

\begin{proof}
It is enough to treat cyclic groups and then use direct-sum
factorization.  For $\ZZ$, one has
\[
\Sigma_{\ZZ}(N)=|\mu_N|=N=|\ZZ/N\ZZ|.
\]
For $C_m=\ZZ/m\ZZ$, a homomorphism $C_m\to\mu_N$ is determined by the
image of a generator, and this image must be an element of $\mu_N$ of
order dividing~$m$.  Hence
\[
\Sigma_{C_m}(N)=\gcd(m,N).
\]
On the other hand,
\[
|C_m/NC_m|=\gcd(m,N).
\]
Taking products over the cyclic decomposition of~$G$ proves the claim.
\end{proof}

\begin{theorem}[Phase-sampling growth recovers circle dimension]%
\label{thm:cdim-phase-spectrum-limit}
Let $G$ be a finitely generated abelian group.  Then
\[
\lim_{N\to\infty}\delta_G(N)=\cdim(G).
\]
More precisely, if $p$ is prime and $a\to\infty$, then
\[
\delta_G(p^a)=\cdim(G)+O(a^{-1}).
\]
\end{theorem}

\begin{proof}
Write $G\simeq \ZZ^r\oplus T$ with $T$ finite.  Then
\[
\Sigma_G(N)=\Sigma_{\ZZ^r}(N)\Sigma_T(N)=N^r\Sigma_T(N).
\]
Moreover
\[
1\le \Sigma_T(N)\le |T|.
\]
Therefore
\[
\delta_G(N)=r+\frac{\log\Sigma_T(N)}{\log N}\longrightarrow r.
\]
By Proposition~\ref{prop:circle-dimension-basic}, $r=\cdim(G)$.
For $N=p^a$ one obtains the quantitative bound
\[
\bigl|\delta_G(p^a)-\cdim(G)\bigr|
\le \frac{\log |T|}{a\log p}.
\]
\end{proof}

\begin{proposition}[Phase-coordinate interface]%
\label{prop:phase-coordinate-interface}
Let
\[
z(\theta):=\iota(\tan\theta)=e^{2\ii\theta}
\qquad\bigl(-\pi/2<\theta<\pi/2\bigr).
\]
Then the analytic phase modes used in
Theorem~\ref{thm:cayley-chebyshev-mode} are the real and imaginary parts
of the circle characters:
\[
\cos(2k\theta)=\operatorname{Re} z(\theta)^k,
\qquad
\sin(2k\theta)=\operatorname{Im} z(\theta)^k
\qquad(k\ge1).
\]
For every finitely generated abelian group~$G$ and every $N\ge1$,
\[
\operatorname{Hom}(G,\mu_N)
=\widehat G[N]
:=\{\chi\in\widehat G:\chi^N=1\},
\]
where $\chi^N(g)=\chi(g)^N$.  Hence
\[
\Sigma_G(N)=|\widehat G[N]|.
\]
Finally, if $K<N$, then normalized averaging over the same phase grid
$\mu_N$ agrees with Haar averaging on every trigonometric monomial of
Fourier degree at most~$K$:
\[
\frac1N\sum_{\zeta\in\mu_N}\zeta^\ell
=\int_{\TT} z^\ell\,\dd m_{\TT}(z)
=
\begin{cases}
1,& \ell=0,\\
0,& 0<|\ell|\le K.
\end{cases}
\]
Consequently, the finite phase averages used in the entropy section are
ordinary circle-character averages for the same Cayley phase coordinate
whose analytic Fourier characters drive the entropy and RKHS arguments
below.
\end{proposition}

\begin{proof}
The identity $z(\theta)=e^{2\ii\theta}$ follows immediately from
\eqref{eq:cayley-embed} with $x=\tan\theta$.  Taking real and imaginary
parts of $z(\theta)^k=e^{2k\ii\theta}$ gives the displayed formulas for
$\cos(2k\theta)$ and $\sin(2k\theta)$.

A character $\chi\in\widehat G=\operatorname{Hom}(G,\TT)$ satisfies
$\chi^N=1$ exactly when $\chi(g)^N=1$ for every $g\in G$, equivalently
when the image of~$\chi$ is contained in~$\mu_N$.  This proves
$\operatorname{Hom}(G,\mu_N)=\widehat G[N]$ and hence
$\Sigma_G(N)=|\widehat G[N]|$ by Definition~\ref{def:cdim-phase-spectrum}.

For the grid average, the geometric sum gives
\[
\frac1N\sum_{\zeta\in\mu_N}\zeta^\ell
=
\begin{cases}
1,& N\mid \ell,\\
0,& N\nmid \ell.
\end{cases}
\]
If $0<|\ell|\le K<N$, then $N\nmid\ell$, while the case $\ell=0$ gives
one.  The Haar integral of the character $z^\ell$ has the same value for
these $\ell$.  Since any monomial $z^m\overline z^{\,n}$ equals
$z^{m-n}$, the assertion covers all monomials with
$0\le m,n\le K$.
\end{proof}

\begin{proposition}[Finite-grid coefficient recovery and aliasing]%
\label{prop:finite-grid-coefficient-recovery}
Let $f$ be a continuous function on~$\TT$ with absolutely convergent
Fourier series
\[
f(z)=\sum_{m\in\ZZ}\widehat f(m)z^m,
\qquad
\sum_{m\in\ZZ}|\widehat f(m)|<\infty,
\]
where
\[
\widehat f(m):=\int_{\TT}f(z)z^{-m}\,\dd m_{\TT}(z).
\]
For $N\ge1$ and $k\in\ZZ$ define the finite phase-grid coefficient recovery
\[
\widehat f_N(k):=\frac1N\sum_{\zeta\in\mu_N}f(\zeta)\zeta^{-k}.
\]
Then
\begin{equation}\label{eq:finite-grid-aliasing}
\widehat f_N(k)=\sum_{q\in\ZZ}\widehat f(k+qN),
\end{equation}
and hence
\begin{equation}\label{eq:finite-grid-aliasing-bound}
\bigl|\widehat f_N(k)-\widehat f(k)\bigr|
\le
\sum_{\substack{q\in\ZZ\\q\ne0}}\bigl|\widehat f(k+qN)\bigr|.
\end{equation}
In particular, if $K<N$ and
\[
\widehat f(m)=0\qquad\text{whenever } |m-k|>K,
\]
then the Fourier support of $z^{-k}f(z)$ lies below the grid scale and the
recovery is exact:
\[
\widehat f_N(k)=\widehat f(k).
\]
For a real phase profile $F(\theta)=f(e^{2\ii\theta})$ on
$(-\pi/2,\pi/2)$, put
\[
\mathcal C_k(F):=2\operatorname{Re}\widehat f(k),
\qquad
\mathcal S_k(F):=-2\operatorname{Im}\widehat f(k),
\]
and define $\mathcal C_{k,N}(F),\mathcal S_{k,N}(F)$ by the same formulas
with $\widehat f_N(k)$ in place of $\widehat f(k)$.  Then
\[
|\mathcal C_{k,N}(F)-\mathcal C_k(F)|
+|\mathcal S_{k,N}(F)-\mathcal S_k(F)|
\le
4\sum_{\substack{q\in\ZZ\\q\ne0}}\bigl|\widehat f(k+qN)\bigr|.
\]
Thus finite phase grids recover the phase-Fourier coefficients used later
in Proposition~\ref{prop:cdim-poisson-cauchy-moment-inversion-by-phase-modes}
exactly for finite mode polynomials below the explicit $K<N$ barrier, and
with only the displayed aliasing tail for general absolutely summable
profiles.
\end{proposition}

\begin{proof}
The absolute convergence assumption permits termwise averaging:
\[
\widehat f_N(k)
=\sum_{m\in\ZZ}\widehat f(m)
\left(\frac1N\sum_{\zeta\in\mu_N}\zeta^{m-k}\right).
\]
The inner average is $1$ when $m-k$ is divisible by~$N$ and is $0$
otherwise.  This proves \eqref{eq:finite-grid-aliasing}.  Removing the
$q=0$ term and applying the triangle inequality gives
\eqref{eq:finite-grid-aliasing-bound}.

If the Fourier support of $z^{-k}f$ is contained in $\{-K,\dots,K\}$, then
$\widehat f(k+qN)=0$ for every $q\ne0$, because
$|k+qN-k|=|q|N\ge N>K$.  The exact recovery follows.

Finally,
\[
\int_{-\pi/2}^{\pi/2}F(\theta)e^{-2\ii k\theta}\,\frac{\dd\theta}{\pi}
=\widehat f(k),
\]
so for real $F$ the cosine and sine coefficients are
$2\operatorname{Re}\widehat f(k)$ and $-2\operatorname{Im}\widehat f(k)$.
The same identities with the grid coefficient give the finite recoveries, and
the two displayed coefficient bounds follow from
\eqref{eq:finite-grid-aliasing-bound}.
\end{proof}
